\documentclass[a4paper,twocolumn,landscape]{article}
\usepackage{geometry}
\usepackage{concmath}
\usepackage[OT1]{fontenc}
\usepackage{graphicx}
\usepackage{amsfonts}
\usepackage{amsmath}
\usepackage{calrsfs}
\usepackage{color}
\usepackage[dvipsnames]{xcolor}
\usepackage[many]{tcolorbox}
\usepackage{setspace}              
\usepackage{multicol}  
\geometry{margin=0.5in}
\definecolor{green}{RGB}{52, 235, 168}
\definecolor{violet}{RGB}{165, 83, 237}
\definecolor{lemon}{RGB}{170, 237, 83}
\definecolor{red}{RGB}{240, 125, 84}
\definecolor{blue}{RGB}{105, 255, 242}

\begin{document}
\noindent
2/ Supposons P(L) (\(\tilde{L}\) régulier), mq P(L*) (\(\tilde{L*}\) régulier)\\
\begin{equation*}
    \tilde{L*} = \tilde{L}*
\end{equation*}
\begin{align*}
    u \in \tilde{L*} &\Leftrightarrow \exists u_1,...,u_n \in L, u = \tilde{u_1...u_n}\\
    &\Leftrightarrow ..., u = \tilde{u_n}...\tilde{u_1}\\
    &\Leftrightarrow u \in \tilde{L}*
\end{align*}
Donc \(\tilde{L*}\) régulier\\

\noindent \colorbox{red}{Complément : Dénombrabilité et langages non réguliers}\\
\colorbox{blue}{Théorème de Cantor} \\
Si \(\exists x, f(x) = Y\):\\
- Si \(x \in f(x): x\notin Y\) absurde car \(f(x) = Y\)\\
- Si \(x \notin f(x): x \in Y\) absurde car \(f(x) = Y\)\\
Donc f n'est pas surjective

\noindent
Soit A un automate déterministe complet, q un état et u, v des mots.
Alors :
\begin{equation*}
    \delta*(q, uv) = \delta*(\delta *(q, u), v)
\end{equation*}
\(v \in \Sigma^*\)\\
H(n) : "Soient q état, si \(|u| = n\) alors \(\delta^*(q,uv) = \delta^*(\delta^*(q,u),v)\)"\\
* H(0) vraie car \(\delta^*(q,\varepsilon)\) = q\\
* Supposons H(n). Soit \(u \in \Sigma*, q\in Q\) tq \(|u| = n+1\)\\
u = au' avec \(a \in \Sigma, |u'| = n\)\\
\begin{align*}
    \delta*(q,uv) &= \delta*(\delta(q,a), u'v) = \delta*(\delta*(\delta(q,a),u'),v)\\
    \text{ avec }\delta*(\delta(q,a),u') &= \delta*(q,au')
\end{align*}
Donc H(n+1)


\end{document}